\renewcommand{\thetheorem}{9.\arabic{theorem}}
\setcounter{theorem}{0}
\setcounter{chapter}{8}

\chapter{Linear Independence}

\subsection{Linear Independence of a Finite List}

Let $V$ be a vector space over a field $K$. In the previous chapter, we studied how a set of vectors can \textit{span} a subspace. We now ask whether a spanning set contains "redundant" vectors. This leads us to the core algebraic concept of linear independence.

\begin{definition}[Linear Independence]
Let $v_1, \dots, v_n$ be a finite list of vectors in $V$. We say that the list $v_1, \dots, v_n$ is \textbf{linearly independent} if the only linear combination that equals the zero vector $0_V$ is the trivial combination (where all coefficients are zero). 

In formal notation, the vectors are linearly independent if, for all $a_1, \dots, a_n \in K$:
\[
a_1 v_1 + a_2 v_2 + \dots + a_n v_n = 0_V \implies a_1 = 0, a_2 = 0, \dots, a_n = 0
\]
\end{definition}

\begin{remark}[The Empty Set]
By mathematical convention, the empty list of vectors $\emptyset$ is defined to be linearly independent. But \textit{why} does this make sense? 

Recall from Chapter 8 that we defined the span of the empty set as the zero subspace: $\Sp(\emptyset) = \{0_V\}$. If the empty set were to be linearly \textit{dependent}, there would need to exist a linear combination of its elements that non-trivially yields $0_V$. However, since $\emptyset$ contains literally no vectors, we cannot even form a linear combination with a non-zero coefficient! Because the condition for dependence is impossible to satisfy, the empty set is vacuously, and logically, independent.
\end{remark}

Before looking at examples, we can immediately deduce several structural facts directly from the definition.

\begin{lemma}[Basic Properties of Linear Independence]
Let $V$ be a vector space. The following properties hold for any finite list of vectors:
\begin{enumerate}[label=\textbf{(\alph*)}]
    \item The zero vector $0_V$ can always be written as a trivial linear combination of any finite list of vectors: $0 \cdot v_1 + \dots + 0 \cdot v_n = 0_V$.
    \item The order of the vectors in a list does not affect their linear independence.
    \item If a list $v_1, \dots, v_n$ contains two identical vectors (i.e., $v_k = v_l$ for some $k \ne l$), then the list is NOT linearly independent.
    \item If the zero vector $0_V$ is an element of the list, then the list is NOT linearly independent.
\end{enumerate}
\end{lemma}

\begin{proof}
We address each property:
\begin{enumerate}[label=\textbf{(\alph*)}]
    \item Follows trivially from the vector space axioms (specifically $0 \cdot v = 0_V$ for all $v \in V$).
    \item Follows trivially from the commutativity of vector addition in $V$.
    \item Assume $v_k = v_l$ for some $1 \le k < l \le n$. We can construct a specific linear combination where the coefficient for $v_k$ is $1$, the coefficient for $v_l$ is $-1$, and all other coefficients are $0$:
    \[
    0 \cdot v_1 + \dots + 1 \cdot v_k + \dots + (-1) \cdot v_l + \dots + 0 \cdot v_n = v_k - v_l = 0_V
    \]
    Because we found a linear combination equal to $0_V$ where not all coefficients are zero (specifically, $a_k = 1 \ne 0$), the list fails the definition and is therefore not linearly independent.

    \item Suppose $v_1 = 0_V$. We can assign a non-zero coefficient (e.g., $1$) to $v_1$ and zero to all other vectors:
    \[
    1 \cdot 0_V + 0 \cdot v_2 + \dots + 0 \cdot v_n = 0_V + 0_V + \dots + 0_V = 0_V
    \]
    Again, since $a_1 = 1 \ne 0$, we have a non-trivial combination yielding the zero vector. The list is not linearly independent.
\end{enumerate}
\end{proof}

\subsection{Examples of Linear Independence}

Let us apply this definition to various vector spaces.

\begin{example}[In $\mathbb{R}^2$]
Consider the vectors $v_1 = \binom{1}{0}$ and $v_2 = \binom{0}{1}$ in $\mathbb{R}^2$. \\
Assume $a_1 v_1 + a_2 v_2 = 0_{\mathbb{R}^2}$. This means:
\[
a_1 \begin{pmatrix} 1 \\ 0 \end{pmatrix} + a_2 \begin{pmatrix} 0 \\ 1 \end{pmatrix}
    = \begin{pmatrix} a_1 \\ a_2 \end{pmatrix}
    = \begin{pmatrix} 0 \\ 0 \end{pmatrix}.
\]
Thus, $a_1 = 0$ and $a_2 = 0$. The vectors $\{v_1, v_2\}$ are linearly independent.
\end{example}

\begin{example}[Linearly dependent vectors in $\mathbb{R}^2$]
Consider $v_1 = \binom{1}{1}$ and $v_2 = \binom{2}{2}$. \\
We can easily spot that $2 v_1 - v_2 = 0_{\mathbb{R}^2}$. Because we have coefficients $a_1 = 2$ and $a_2 = -1$ (which are non-zero) that yield the zero vector, these vectors are NOT linearly independent.
\end{example}

\begin{example}[The Standard Basis in $K^n$]
Consider the standard basis vectors in $K^n$:
\[
e_1 = \begin{pmatrix} 1 \\ 0 \\ \vdots \\ 0 \end{pmatrix}, \quad e_2 = \begin{pmatrix} 0 \\ 1 \\ \vdots \\ 0 \end{pmatrix}, \quad \dots \quad, \quad e_n = \begin{pmatrix} 0 \\ 0 \\ \vdots \\ 1 \end{pmatrix}
\]
Setting $a_1 e_1 + a_2 e_2 + \dots + a_n e_n = 0_{K^n}$ yields the column vector $(a_1, a_2, \dots, a_n)\transp = (0, 0, \dots, 0)\transp$. This immediately forces $a_i = 0$ for all $i$. Thus, the set $\{e_1, e_2, \dots, e_n\}$ is always linearly independent.
\end{example}

\begin{example}[Single Vectors]
Let $V$ be a vector space and $v \in V$. The single-element set $\{v\}$ is linearly independent if and only if $v \ne 0_V$. (If $v = 0_V$, we know from Lemma 9.1 that it is not independent. If $v \ne 0_V$, then $a \cdot v = 0_V \implies a = 0$).
\end{example}

\subsection{Arbitrary Families and Linear Dependence}

As we advance in linear algebra, we frequently need to deal with infinite collections of vectors (such as the set of all polynomials). We must rigorously extend our definition to handle this.

\begin{definition}[Linear Independence of an Arbitrary Family]
Let $\mathcal{F} = \{v_i\}_{i \in I}$ be a \textit{family of vectors} in $V$, indexed by an arbitrary set $I$ (which may be infinite). We say that the family $\mathcal{F}$ is linearly independent if \textbf{every finite subset} of $\mathcal{F}$ is linearly independent.
\end{definition}

\begin{example}[Families of Polynomials]
In the space of bounded polynomials $K[x]_d$, the set of standard monomials $\{1, x, x^2, \dots, x^d\}$ is linearly independent. This remains true even for the \textbf{infinite family} of all monomials $\mathcal{F} = \{1, x, x^2, \dots\} \subseteq K[x]$. Any finite subset of this family will just be a finite collection of distinct monomials, which can only sum to the zero polynomial if all their coefficients are strictly zero.
\end{example}

We now formally define the opposite of linear independence, specifically focusing on how it applies to an arbitrary family.

\begin{definition}[Linear Dependence]
A family $\mathcal{F} = \{v_i\}_{i \in I}$ of vectors in $V$ is called \textbf{linearly dependent} if it is NOT linearly independent. 

In other words, this means there exists some finite integer $n \in \mathbb{N}$, a selection of $n$ distinct indices $i_1, \dots, i_n \in I$, and a set of scalars $a_1, \dots, a_n \in K$ \textbf{where not all $a_j$ are zero}, such that:
\[
a_1 v_{i_1} + a_2 v_{i_2} + \dots + a_n v_{i_n} = 0_V
\]
\end{definition}

\subsection{Back to Linear Independence: Redundancy}

This definition leads to an incredibly useful geometric intuition: a set of vectors is linearly dependent if and only if at least one vector is "redundant" — meaning it can be built entirely out of the other vectors in the set. 

\begin{center}
\begin{tikzpicture}[
    node distance=2cm,
    box/.style={draw, thick, rounded corners, fill=blue!5, align=center, inner sep=2ex, text width=5cm},
    arrow/.style={<->, thick, >=stealth, draw=blue!60!black}
]
    \node[box] (dep) {\textbf{Algebraic Definition} \\ $\{v_1, \dots, v_n\}$ is \\ Linearly Dependent \\ ($\sum a_i v_i = 0_V$, $a_i \ne 0$)};
    \node[box, right=3cm of dep] (red) {\textbf{Geometric Redundancy} \\ At least one vector $v_k$ \\ is a linear combination \\ of the others. \\ ($v_k \in \Sp(v_1, \dots, \widehat{v_k}, \dots, v_n)$)};
    
    \draw[arrow] (dep) -- (red) node[midway, above, font=\bfseries] {Proposition 9.1};
\end{tikzpicture}
\end{center}

\begin{proposition}[Equivalence of Dependence and Redundancy]
A list of vectors $v_1, \dots, v_n$ is linearly dependent if and only if at least one vector in the list can be written as a linear combination of the other vectors.
\end{proposition}

\begin{proof}
We must prove both directions.

\textbf{Direction 1 ($\implies$):} 
Assume the list $v_1, \dots, v_n$ is linearly dependent. By definition, there exist scalars $a_1, \dots, a_n \in K$, not all zero, such that:
\[
a_1 v_1 + a_2 v_2 + \dots + a_n v_n = 0_V
\]
Since not all coefficients are zero, we can choose an index $k$ such that $a_k \ne 0$. We can isolate $a_k v_k$ on one side of the equation by subtracting all other terms:
\[
a_k v_k = - \sum_{\substack{i=1 \\ i \ne k}}^n a_i v_i
\]
Because $a_k \ne 0$, its multiplicative inverse $a_k^{-1}$ exists in the field $K$. Multiplying both sides by $a_k^{-1}$ yields:
\[
v_k = \sum_{\substack{i=1 \\ i \ne k}}^n \left( -\frac{a_i}{a_k} \right) v_i
\]
Thus, $v_k$ is a linear combination of the other vectors.

\textbf{Direction 2 ($\Longleftarrow$):} 
Assume that at least one vector, say $v_k$, can be written as a linear combination of the others. This means there exist scalars $b_i$ such that:
\[
v_k = \sum_{\substack{i=1 \\ i \ne k}}^n b_i v_i
\]
We can subtract $v_k$ from both sides to obtain:
\[
0_V = -1 \cdot v_k + \sum_{\substack{i=1 \\ i \ne k}}^n b_i v_i
\]
This equation is a linear combination of the full list $v_1, \dots, v_n$ that equals $0_V$. Furthermore, the coefficient of $v_k$ is $-1$, which is non-zero. Since we have found a non-trivial linear combination equaling $0_V$, the list is linearly dependent by definition.
\end{proof}

\subsection{Inheritance in Subsets and Supersets}

How does linear independence behave when we add or remove vectors from a set? There is a strict, logical hierarchy to how these properties are inherited. 

\begin{center}
\begin{tikzpicture}[
    box/.style={draw, thick, rounded corners, align=center, inner sep=2ex},
    arrow/.style={->, thick, >=stealth}
]
    % Left Diagram: Independence goes down
    \begin{scope}[shift={(0,0)}]
        \node (center1) at (0,0) {}; % Invisible anchor node
        \draw[thick, fill=green!10] (center1) ellipse (2.5cm and 1.8cm);
        \node[above] at (0, 1) {\textbf{Set $S$}};
        \node[above] at (0, 0.5) {(Linearly Independent)};
        
        \draw[thick, fill=green!30, dashed] (0,-0.6) ellipse (1.2cm and 0.8cm);
        \node at (0, -0.6) {\textbf{Subset $S'$}};
        
        \node[box, fill=white, below=2.2cm of center1] (l1) {\textbf{Lemma 9.2} \\ Independence is \\ inherited downwards.};
        \draw[arrow, color=green!50!black] (0, 0.2) -- (0, -0.2);
    \end{scope}

    % Right Diagram: Dependence goes up
    \begin{scope}[shift={(7,0)}]
        \node (center2) at (0,0) {}; % Invisible anchor node
        \draw[thick, fill=red!10] (center2) ellipse (2.5cm and 1.8cm);
        \node[above] at (0, 1) {\textbf{Superset $S''$}};
        
        \draw[thick, fill=red!30, dashed] (0,-0.6) ellipse (1.2cm and 0.8cm);
        \node at (0, -0.6) {\textbf{Set $S$}};
        \node at (0, -1) {(Dependent)};
        
        \node[box, fill=white, below=2.2cm of center2] (l2) {\textbf{Lemma 9.3} \\ Dependence \\ spreads upwards.};
        \draw[arrow, color=red!50!black] (0, -0.2) -- (0, 0.6);
    \end{scope}
\end{tikzpicture}
\end{center}

\begin{lemma}[Subsets of Independent Sets]
Let $S \subseteq V$ be a linearly independent set. If $S' \subseteq S$ is a subset, then $S'$ is also linearly independent.
\end{lemma}

\begin{proof}
We prove this by contradiction. Suppose $S'$ is NOT linearly independent. Then $S'$ must be linearly dependent. By definition, there exist vectors $v_1, \dots, v_k \in S'$ and scalars $a_1, \dots, a_k \in K$, not all zero, such that:
\[
a_1 v_1 + \dots + a_k v_k = 0_V
\]
However, since $S' \subseteq S$, every vector $v_i$ is also an element of the larger set $S$. Therefore, this exact same equation represents a non-trivial linear combination of elements in $S$ that equals $0_V$. This would mean $S$ is linearly dependent, which directly contradicts our initial assumption that $S$ is linearly independent. Thus, $S'$ must be linearly independent.
\end{proof}

\begin{lemma}[Supersets of Dependent Sets]
Let $S \subseteq V$ be a linearly dependent set. If $S''$ is a superset containing $S$ (i.e., $S \subseteq S'' \subseteq V$), then $S''$ is also linearly dependent.
\end{lemma}

\begin{proof}
Because $S$ is linearly dependent, there exist vectors $v_1, \dots, v_k \in S$ and scalars $a_1, \dots, a_k \in K$, not all zero, such that:
\[
a_1 v_1 + \dots + a_k v_k = 0_V
\]
Since $S \subseteq S''$, the vectors $v_1, \dots, v_k$ also belong to $S''$. Therefore, this exact same non-trivial linear combination exists within $S''$, immediately satisfying the definition of linear dependence for $S''$. 
\end{proof}
